\begin{definition}[Carathéodory Factorization]
Let $I \subseteq \mathbb{R}$ be an interval, let $f : I \to \mathbb{R}$, and let $c \in I$.
We say that $f$ admits a \emph{Carathéodory factorization} at $c$ if there exists a function
$\phi : I \to \mathbb{R}$, continuous at $c$, such that
\[
f(x)=f(c)+(x-c)\phi(x)
\qquad \text{for all } x \in I.
\]
In this case, $f$ is differentiable at $c$, and
\[
f'(c)=\phi(c).
\]
\end{definition}


\begin{theorem}[Carath\'eodory Characterization of Differentiability]
\label{thm:caratheodory-characterization-of-differentiability}
Let $I \subseteq \mathbb{R}$ be an interval, let $f : I \to \mathbb{R}$, and let $c \in I$. Then the following are equivalent:

\begin{enumerate}[label=(\roman*)]
\item $f$ is differentiable at $c$.

\item There exists a function $\phi : I \to \mathbb{R}$, continuous at $c$, such that
\[
f(x)=f(c)+(x-c)\phi(x)
\qquad \text{for all } x \in I.
\]
\end{enumerate}

In this case,
\[
\phi(c)=f'(c).
\]
\end{theorem}


\begin{remark}[The explicit factorizing function]
The unique $\varphi$ satisfying the identity is
\[
\varphi(x) =
\begin{cases}
\dfrac{f(x) - f(c)}{x - c} & x \neq c, \\[6pt]
L & x = c.
\end{cases}
\]
Differentiability of $f$ at $c$ is precisely the statement that this
piecewise-defined function is continuous at $c$.
The factorization eliminates the division-by-zero that makes the
limit definition awkward in algebraic arguments.
\end{remark}

\begin{remark}[Chain rule via Carathéodory]
If $f$ is differentiable at $c$ with factor $\varphi$, and $g$ is
differentiable at $f(c)$ with factor $\psi$, then
\begin{align*}
g(f(x)) - g(f(c))
  &= \psi(f(x))\cdot(f(x) - f(c)) \\
  &= \psi(f(x))\cdot\varphi(x)\cdot(x - c).
\end{align*}
The function $x \mapsto \psi(f(x))\cdot\varphi(x)$ is continuous at $c$
and evaluates there to $g'(f(c))\cdot f'(c)$.
This gives the chain rule with no limit computation.
\end{remark}

\begin{remark}[Why this matters]
The Carathéodory form is the preferred tool for proving the product rule,
quotient rule, and chain rule cleanly. It converts limit arguments into
continuity arguments, which are often simpler to chain together.
\end{remark}






\begin{remark*}
Differentiability at $c$ is equivalent to the ability to factor the increment
$f(x) - f(c)$ by $(x-c)$ with a coefficient function $\varphi(x)$ that extends
continuously to $x = c$. The derivative is then simply recovered as
$f'(c) = \varphi(c)$. This reframes a limit condition on a quotient as a
continuity condition on a function — a reformulation that makes algebraic
manipulation of derivatives significantly cleaner.
\end{remark*}

\begin{remark*}[Carathéodory factorization vs.\ Lipschitz control]
Both Carathéodory's theorem and the Lipschitz condition control the increment
$f(x) - f(c)$ by the factor $(x-c)$, but in fundamentally different ways.

\textbf{Carathéodory factorization} gives an exact linear factorization: the
slope function $(f(x)-f(c))/(x-c)$ is required to extend continuously to
$x = c$. The derivative appears as the value of this continuous extension.
This is a differentiability criterion — it controls the behaviour of the slope,
not merely its size.

\textbf{Lipschitz control} gives only an inequality: $|f(x)-f(c)| \le L|x-c|$
for $x$ near $c$. This bounds how large the increment can be, but imposes no
requirement that the slope quotient converge to any limit.

\textbf{Relationship.} Carathéodory's representation implies Lipschitz: if
$\varphi$ is continuous at $c$, it is bounded near $c$, so
$|f(x)-f(c)| = |\varphi(x)||x-c| \le L|x-c|$ for some $L$ and all $x$ near
$c$. Thus differentiability at $c$ implies a local Lipschitz estimate. The
converse fails: $f(x) = |x|$ is Lipschitz on $\mathbb{R}$ with constant $L = 1$,
but the slope function takes values $\pm 1$ on either side of $0$ with no
continuous extension, so $f$ is not differentiable at $0$.
\end{remark*}




\begin{theorem}[Chain Rule]
Let $I, J \subseteq \mathbb{R}$ be intervals, let $f : I \to \mathbb{R}$ and $g : J \to \mathbb{R}$, and suppose that $f(I) \subseteq J$. Let $c \in I$.

If $f$ is differentiable at $c$ and $g$ is differentiable at $f(c)$, then the composition $g \circ f$ is differentiable at $c$, and
\[
(g \circ f)'(c) = g'(f(c))\,f'(c).
\]
\end{theorem}







\begin{remark*}
The standard approach via Carathéodory's theorem avoids the division-by-zero
failure of the naive proof. Let $\varphi_f$ be the Carathéodory function of
$f$ at $c$ and $\varphi_g$ be the Carathéodory function of $g$ at $f(c)$.
Then:
\[
  g(f(x)) - g(f(c))
  = \varphi_g(f(x)) \cdot (f(x) - f(c))
  = \varphi_g(f(x)) \cdot \varphi_f(x) \cdot (x - c).
\]
The function $x \mapsto \varphi_g(f(x)) \cdot \varphi_f(x)$ is continuous at
$c$: $\varphi_f$ is continuous at $c$ by differentiability of $f$, $f$ is
continuous at $c$ (differentiability implies continuity), and $\varphi_g$ is
continuous at $f(c)$ by differentiability of $g$. The value at $x = c$ is
$\varphi_g(f(c)) \cdot \varphi_f(c) = g'(f(c)) \cdot f'(c)$. By
Carathéodory's theorem applied to $g \circ f$, this gives the result.
\end{remark*}

